%%%%%%%%%%%%%%%%%%%%%%%%%%%%%%%%%%%%%%%%%%%%%%%%%%%%%%%%%%%%%%%%%%%%%%
% Lớp 8 | Chương I. Đa thức | Phép nhân đa thức
% Dạng 8D1?4-4: Tìm x, bài toán thực tế (diện tích, thể tích)
%%%%%%%%%%%%%%%%%%%%%%%%%%%%%%%%%%%%%%%%%%%%%%%%%%%%%%%%%%%%%%%%%%%%%%


% nguồn: De_AI_Toan8_cuoi_hoc_ki_i_1 (nhập 21/11/2026)
% nguồn: De_AI_Toan8_cuoi_hoc_ki_i_1 (nhập 21/11/2026)
% nguồn: De_AI_Toan8_cuoi_hoc_ki_i_1 (nhập 21/11/2026)
% nguồn: De_AI_Toan8_cuoi_hoc_ki_i_1 (nhập 21/11/2026)
% nguồn: De_AI_Toan8_cuoi_hoc_ki_i_1 (nhập 21/11/2026)
\begin{bt}%[8D1H4-4]%[ma:4921bb1a]
	Tìm ba số tự nhiên liên tiếp, biết rằng tích của hai số sau lớn hơn tích của hai số đầu là $24$.
	\loigiai{
		Gọi ba số tự nhiên liên tiếp là $n$, $n + 1$, $n + 2$ ($n \in \mathbb{N}$). \\
Theo đề bài ta có $(n + 1)(n + 2) - n(n + 1) = 24$. \\
$(n^2 + 3n + 2) - (n^2 + n) = 24 \Rightarrow 2n + 2 = 24 \Rightarrow 2n = 22 \Rightarrow n = 11$. \\
Vậy ba số tự nhiên liên tiếp cần tìm là $11$, $12$, $13$.
	}
\end{bt}

% nguồn: De_AI_Toan8_cuoi_hoc_ki_i_1 (nhập 21/11/2026)
\begin{bt}%[8D1H4-4]%[ma:afa5ca81]
	Một bể cá dạng hình hộp chữ nhật có ba kích thước là $x$, $x + 2$ và $x + 3$ (dm) với $x > 0$.
a) Viết đa thức thu gọn biểu thị thể tích của bể cá.
b) Tính thể tích bể cá (theo đơn vị lít, biết $1\text{ dm}^3 = 1\text{ lít}$) khi $x = 4$.
	\loigiai{
		a) Thể tích bể cá là $V = x(x + 2)(x + 3) = x(x^2 + 5x + 6) = x^3 + 5x^2 + 6x$ ($\text{dm}^3$). \\
b) Thay $x = 4$ vào $V$, ta được $V = 4 \cdot 6 \cdot 7 = 168\text{ dm}^3 = 168\text{ lít}$.
	}
\end{bt}
